\documentclass[a4paper,10pt]{amsart}

\pdfinfoomitdate=1
\pdftrailerid{}
\pdfsuppressptexinfo=15

\usepackage[T1]{fontenc}
\usepackage{lmodern}
\usepackage[margin=28mm]{geometry}
\usepackage{amsmath,amssymb,mathtools,booktabs,microtype,xcolor}
\usepackage[numbers,sort&compress]{natbib}
\usepackage[colorlinks=true,linkcolor=blue!55!black,citecolor=blue!55!black,urlcolor=blue!55!black]{hyperref}
\hypersetup{pdftitle={},pdfauthor={},pdfsubject={},pdfkeywords={}}

\newtheorem{theorem}{Theorem}
\newtheorem{proposition}[theorem]{Proposition}
\newtheorem{lemma}[theorem]{Lemma}
\theoremstyle{remark}
\newtheorem{remark}[theorem]{Remark}

\newcommand{\Prob}{\mathbb P}
\newcommand{\cT}{\mathcal T}

\title[Uniform ear deletion]{Uniform Ear Deletion of a Convex Polygon:\\
Exact Triangulation Masses and a Sharp Path-Dual Minimum}
\author{Anonymous}
\date{}

\begin{document}

\begin{abstract}
Start with a labelled convex $n$-gon.  Repeatedly choose a current vertex
uniformly, delete it, and record the chord between its two current
neighbours, stopping at a triangle.  Although all $n!/6$ deletion histories
are equiprobable, their triangulation endpoints need not be; nonuniformity
first occurs at $n=6$.  For a triangulation
$T$ and one of its triangular faces $r$, root the weak dual at $r$ and let
$s_v^{(r)}$ be the rooted-subtree size at $v$.  We prove that the number of
histories with endpoint $T$ and final face $r$ is
\[
  H(T,r)=\frac{(n-3)!}{\prod_{v\ne r}s_v^{(r)}}.
\]
Consequently $\Prob(T)=6H(T)/n!$, where $H(T)=\sum_r H(T,r)$.
Equivalently, $H(T)$ counts complete leaf-deletion orders of the unrooted
weak dual.  This yields the sharp bound $H(T)\ge 2^{n-3}$, with equality
exactly when the weak dual is a path.  An exhaustive exact-arithmetic audit
through $n=9$ checks every endpoint and final-face refinement.  Classical ear
clipping, triangulation enumeration, weak duals, and tree hook formulas are
treated as inputs rather than contributions.
\end{abstract}

\maketitle

\section{The process and the theorem}

Let $P_n$ be a convex polygon whose vertices have distinct labels.  If its
current cyclic vertex set has size at least four, choose one current vertex
uniformly, join its two current neighbours, and delete it.  Stop as soon as
three vertices remain.  The recorded chords form the random endpoint
triangulation $T$.  Ear clipping is classical in polygon triangulation
algorithms~\cite{EderHeldPalfrader2018}, and convex-polygon ear deletion also
appears in bijections for triangulations~\cite{Regev2013}.  Here the object of
study is the endpoint law induced by uniform choice among \emph{all} current
vertices.

For a triangulation $T$, its weak dual $D_T$ has one vertex for each triangular
face and one edge for each pair of faces sharing a diagonal.  It is a tree on
$m=n-2$ vertices.  If a face $r$ is distinguished and $D_T$ is rooted at $r$,
write $s_v^{(r)}$ for the number of vertices in the subtree rooted at $v$.
Define
\begin{equation}\label{eq:root-hook}
 H(T,r):=\frac{(n-3)!}{\prod_{v\in V(D_T)\setminus\{r\}}s_v^{(r)}},
 \qquad H(T):=\sum_{r\in V(D_T)}H(T,r).
\end{equation}

\begin{theorem}[complete endpoint law and sharp minimum]\label{thm:main}
For every labelled convex $n$-gon, $n\ge3$, the following hold.
\begin{enumerate}
\item The process lasts $n-3$ steps, and each complete deletion history has
probability $6/n!$.
\item For every triangulation $T$ and face $r$, exactly $H(T,r)$ histories
end at $T$ with $r$ as the final face.  In particular,
\begin{equation}\label{eq:endpoint-law}
       \Prob(T)=\frac{6}{n!}H(T).
\end{equation}
\item The endpoint multiplicity obeys
\begin{equation}\label{eq:lower}
       H(T)\ge 2^{n-3}.
\end{equation}
Equality holds if and only if $D_T$ is a path.  Thus the least endpoint
probability is $6\,2^{n-3}/n!$, attained exactly by path-dual triangulations.
\end{enumerate}
\end{theorem}

The theorem gives every endpoint mass, not just the extremum.  Formula
\eqref{eq:root-hook} is integer-valued because it counts histories; it is the
familiar rooted-tree hook expression, used here after identifying the exact
deletion poset.  The forest hook formula and its refinements are classical
\cite{BjornerWachs1989} and receive zero contribution credit.

\section{Histories as dual-tree orders}

At a stage with $k$ current vertices, every vertex of the convex $k$-gon is
an ear.  Hence every ordered list of $n-3$ distinct original labels is a
legal history.  Its probability is
\begin{equation}\label{eq:history-prob}
 \frac1n\frac1{n-1}\cdots\frac14=\frac{3!}{n!},
\end{equation}
and there are $n!/3!$ such lists.  Each deletion inserts the diagonal that
cuts off its ear.  The induction invariant is that the remaining convex
polygon together with the already cut-off ear triangles triangulates the
original polygon.  Thus the next neighbour chord is new and noncrossing.
After $n-3$ steps there are exactly $n-3$ distinct diagonals, hence a
triangulation.

Fix now an endpoint triangulation $T$ and a proposed final face $r$.  Removing
an ear of $T$ removes a leaf of $D_T$; the ear tip uniquely determines the
leaf face, and conversely a leaf face different from $r$ has a unique ear tip
that can be deleted.  Iterating gives the following precise correspondence.

\begin{lemma}[root-resolved bijection]\label{lem:bijection}
Histories ending at $(T,r)$ are in bijection with orders of
$V(D_T)\setminus\{r\}$ in which every child occurs before its parent when
$D_T$ is rooted at $r$.
\end{lemma}

\begin{proof}
Given a history, label each deleted vertex by the triangular face cut off at
that step.  A face can disappear only after all faces below it in the rooted
dual have disappeared, so descendants occur before ancestors.  The last
triangle is $r$.  Conversely, take any such order.  Every prefix removes a
descendant-closed set, so the unremoved faces form the connected
ancestor-closed subtree containing $r$ and triangulate the current convex
polygon.  When a dual vertex $v$ is scheduled, all its children have gone
while its parent remains.  Thus $v$ is a leaf of the remaining dual and its
unique ear tip is currently deletable.
Deleting those tips reconstructs a unique polygon history and records
exactly the diagonals of $T$.  The two constructions are inverse.
\end{proof}

\begin{proposition}[rooted hook count]\label{prop:hook}
The number of orders in Lemma~\ref{lem:bijection} is the quantity
$H(T,r)$ in \eqref{eq:root-hook}.
\end{proposition}

\begin{proof}
Put $q=n-3=m-1$.  For a rooted tree, split the $q$ nonroot vertices into
the branches below the root, of sizes $q_1,\ldots,q_d$.  Orders internal to
different branches may be interleaved in
$\binom{q}{q_1,\ldots,q_d}$ ways.  Applying the same decomposition within
each branch and cancelling the factorials gives
\[
 \frac{q!}{\prod_{v\ne r}s_v^{(r)}}.
\]
The one-vertex case initializes the induction.  Lemma~\ref{lem:bijection}
then gives the asserted history count.
\end{proof}

Summing Proposition~\ref{prop:hook} over the mutually exclusive final faces
gives $H(T)$.  Multiplying by the common history probability
\eqref{eq:history-prob} proves parts (1) and (2) of
Theorem~\ref{thm:main}.  It also proves normalization without a separate
identity: all $n!/6$ histories have exactly one endpoint and final face.

\section{The sharp path-dual minimum}

For an unrooted tree $D$, let $L(D)$ be the number of ways to delete leaves
one at a time until a single vertex remains.  Choosing that final vertex as
$r$ and then orienting all edges toward it partitions these orders by their
surviving vertex.  Lemma~\ref{lem:bijection} therefore gives
\begin{equation}\label{eq:L=H}
       H(T)=L(D_T).
\end{equation}
If $\mathop{Leaf}(D)$ is the leaf set, then
\begin{equation}\label{eq:recurrence}
 L(K_1)=1,
 \qquad
 L(D)=\sum_{v\in\mathop{Leaf}(D)}L(D-v)\qquad (|V(D)|\ge2).
\end{equation}

\begin{proposition}[sharp leaf-order bound]\label{prop:min}
For every tree $D$ on $m\ge1$ vertices,
\[
       L(D)\ge2^{m-1},
\]
with equality if and only if $D$ is a path (including the one-vertex path).
\end{proposition}

\begin{proof}
Proceed by induction on $m$.  The claim is immediate for $m=1$.  Every tree
with $m\ge2$ has at least two leaves, and $D-v$ is a tree on $m-1$ vertices
for each leaf $v$.  By \eqref{eq:recurrence} and induction,
\[
 L(D)\ge |\mathop{Leaf}(D)|\,2^{m-2}\ge2^{m-1}.
\]
If equality holds, then $D$ has exactly two leaves; a finite tree with exactly
two leaves is a path.  Conversely, deleting either endpoint of a path leaves
a path, so equality follows recursively from \eqref{eq:recurrence}.
\end{proof}

Taking $m=n-2$ in Proposition~\ref{prop:min}, using
\eqref{eq:L=H}, proves part (3) of Theorem~\ref{thm:main}.

\section{Exact audit and scope}

An independent standard-library verifier enumerates every ordered deletion
history through $n=9$.  For each endpoint it reconstructs all triangular
faces and the connected weak dual, checks every final-face count against
\eqref{eq:root-hook}, independently computes the leaf-order recurrence,
brute-forces bounded rooted linear extensions, checks total probability in
exact rational arithmetic, and tests the equality class in
Proposition~\ref{prop:min}.  Table
\ref{tab:audit} is frozen output, not proof evidence.

\begin{table}[ht]
\centering
\small
\caption{Exhaustive endpoint audit.  Here $H_{\min}$ and $H_{\max}$ are
endpoint history multiplicities, and the last column counts path-dual
triangulations.}
\label{tab:audit}
\begin{tabular}{@{}rrrrrr@{}}
\toprule
$n$ & histories & triangulations & $H_{\min}$ & $H_{\max}$ & path dual \\
\midrule
3 & 1     & 1   & 1  & 1   & 1   \\
4 & 4     & 2   & 2  & 2   & 2   \\
5 & 20    & 5   & 4  & 4   & 5   \\
6 & 120   & 14  & 8  & 12  & 12  \\
7 & 840   & 42  & 16 & 28  & 28  \\
8 & 6720  & 132 & 32 & 112 & 64  \\
9 & 60480 & 429 & 64 & 316 & 144 \\
\bottomrule
\end{tabular}
\end{table}

The claims are deliberately bounded.  The polygon is convex, selection is
uniform over current vertices, and the theorem does not address nonconvex
visibility, weighted choices, or a classification of the maximum of $H(T)$.
Ear clipping, Catalan enumeration, weak-dual trees, reduction-order/linear-
extension correspondences---including a different-carrier example for
phylogenetic networks \citep{CoronadoPonsRiera2024}---and the generic tree
hook formula receive no priority claim.  A bounded source search did not locate the combined endpoint
law and sharp minimum above, but a search non-hit is not a novelty or priority
certificate.  The residual is elementary and owner-thin.  The manuscript
remains anonymous and under \texttt{HOLD\_EXTERNAL}.

\bibliographystyle{plainnat}
\bibliography{references}

\end{document}
